\SourceLeaf{21}

Die Observatio kann aus den Criterien leicht abgeleitet werden. Z.\,B. war für den cubischen Rest 3
soeben die Bedingung $N\equiv 0\;(\mathrm{mod.}\ 3)$ gewonnen und unter [2] ebenso angegeben. Sind nun $M$ und $N$
gerade, so setze man:
\[ M=2x,\qquad N=6y, \]
so dass aus $4n=M^2+27N^2$ die Darstellbarkeit von $n$ in der quadratischen Form der Determinante $-243$:
\[ n=x^2+243y^2 \]
vermöge ganzzahliger $x, y$ entspringt. Umgekehrt folgt bei der Möglichkeit einer solchen Darstellung von $n$
die Gültigkeit der Congruenz $N\equiv 0\;(\mathrm{mod.}\ 3)$. Sind $M$ und $N$ ungerade, so setze man unter richtiger Auswahl des Vorzeichens:
\[ M\pm\dfrac{N}{3}=4x,\qquad \mp\dfrac{N}{3}=y, \]
woraus man die Darstellung gewinnt:
\[ n=4x^2+2xy+61y^2. \]
Aus dieser Formel würde umgekehrt folgen $N\equiv 0\;(\mathrm{mod.}\ 3)$, so dass 3 stets und nur dann cubischer Rest
von $n$ ist, wenn $n$ durch eine der Formen $(1,0,243)$, $(4,1,61)$, darstellbar ist.

Durch Herrn R. \textsc{Dedekind}, dem ich mehrere für mich sehr förderliche Unterhaltungen über die in
Rede stehenden Gegenstände verdanke, bin ich darauf aufmerksam gemacht, dass für den Fall des cubischen
Nichtrestes 3 von $n$ \textsc{Gauss}' Angaben nicht ganz vollständig sind. Es giebt insgesammt neun Classen ursprünglicher Formen erster Art der Determinante $D=-243$, deren reducirte Formen die folgenden sind:
\[ (1,0,243),\quad (4,\pm1,61),\quad (13,\pm2,19),\quad (7,\pm3,36),\quad (9,\pm3,28). \]
Es fehlen in [1] die Formen $(13,\pm2,19)$. Das Theorem muss demnach so lauten, dass 3 stets und nur
dann cubischer Nichtrest von $n$ ist, wenn $n$ eine Darstellung in einer der Formen gestattet:
\[ n=13xx+4xy+19yy, \]
\[ n=7xx+6xy+36yy, \]
\[ n=9xx+6xy+28yy. \]
Diese Darstellungen erreicht man immer unter richtiger Fixirung des Vorzeichens durch folgende Transformationen:
\[ \text{I. Wenn } M\equiv N\equiv 1\;(\mathrm{mod.}\ 2),\quad M\pm 7N\equiv 0\;(\mathrm{mod.}\ 12), \]
\[ M=5x+7y,\qquad \pm N=x-y. \]
\[ \text{II. Wenn } M\equiv N\equiv 1\;(\mathrm{mod.}\ 2),\quad M\pm 7N\equiv 6\;(\mathrm{mod.}\ 12), \]
\[ M=x+6y,\qquad \mp N=x. \]
\[ \text{III. Wenn } M\equiv N\equiv 0\;(\mathrm{mod.}\ 2),\quad M\pm N\equiv 0\;(\mathrm{mod.}\ 3), \]
\[ M=6x+2y,\qquad \mp N=2y. \]
\par\noindent\hspace*{\parindent}Die über die Zahl 5 unter 1 gemachten Angaben folgen in ähnlicher Weise aus dem Restcriterium
$MN\equiv 0\;(\mathrm{mod.}\ 5)$.

Übrigens sei noch bemerkt, dass sich im Original der Notiz [1] eine bis 223 fortgesetzte Tabelle derjenigen Primzahlen $p=3n+1$ findet, von denen die Zahlen 2, 3, 5, 7 cubische Reste sind.

\hfill \textsc{Fricke.}

\endgroup
